\tcbset{
    breakable=false, enhanced,   
    left*=10pt, right*=10pt,
    top=0pt, bottom=0pt,
    colback=white!10!white,
    colframe=black!75!black,
    fonttitle=\bfseries\large,
    subtitle style={boxrule=0pt,colback=gray!50!white},
}

\lstset{language=python, breaklines=true}

\begin{tcolorbox}[title=Concept score]
\quad \\
\small 
[System] \\
Please act as an impartial judge and evaluate whether the specified concept is 
presented in the sentence fragment provided below. Focus solely on whether the concept is 
clearly incorporated, without regard for grammar or logical coherence. \\ \\
Begin your evaluation by providing a short explanation of whether the concept is 
clearly incorporated. Point out references to the concept in the fragment, noting any 
minimal or awkward incorporations. After providing your explanation, rate the concept's 
relevance on a scale from 0 to 2, where 0 indicates the concept is not present at all, 
1 indicates the concept is somewhat present but minimally or awkwardly incorporated, 
and 2 indicates the concept is more fully and effectively incorporated, with stronger and 
more natural integration. Provide your rating using this exact format: ``Rating: [[score]]''. \\ \\
{[Concept Start]} \\
\textcolor{gray}{[Concept goes here]} \\
{[Concept End]} \\ \\
{[Sentence Fragment Start]} \\
\textcolor{gray}{[Sentence goes here]} \\
{[Sentence Fragment End]} \\
\end{tcolorbox}

\begin{tcolorbox}[title=Instruct score]
\quad \\
\small 
[System] \\
Please act as an impartial judge and evaluate whether the 
sentence fragment provided below is related to the instruction. Focus solely 
on the degree of relatedness in terms of topic, regardless of grammar, coherence, or 
informativeness. \\ \\
Begin your evaluation by providing a brief explanation of whether the 
sentence is related to the instruction, and point out references 
related to the instruction. After providing your explanation, rate the instruction 
relevance on a scale from 0 to 2, where 0 indicates the sentence is unrelated to the 
instruction, 1 indicates it is somewhat related but only minimally or indirectly relevant in terms of topic, 
and 2 indicates it is more clearly and directly related to the instruction. Provide your rating 
using this exact format: ``Rating: [[score]]''. \\ \\
{[Instruction Start]} \\
\textcolor{gray}{[Instruction goes here]} \\ 
{[Instruction End]} \\ \\ 
{[Sentence Fragment Start]} \\ 
\textcolor{gray}{[Sentence goes here]} \\ 
{[Sentence Fragment End]} \\
\end{tcolorbox}

\clearpage

\begin{tcolorbox}[title=Fluency score]
\quad \\
\small 
[System] \\
Please act as an impartial judge and evaluate the fluency of the 
sentence fragment provided below. Focus solely on fluency, disregarding 
its completeness, relevance, coherence with any broader context, or informativeness. \\ \\
Begin your evaluation by briefly describing the fluency of the sentence, noting any 
unnatural phrasing, awkward transitions, grammatical errors, or repetitive structures that 
may hinder readability. After providing your explanation, rate the sentence's fluency 
on a scale from 0 to 2, where 0 indicates the sentence is not fluent and highly unnatural 
(e.g., incomprehensible or repetitive), 1 indicates it is somewhat fluent but contains noticeable 
errors or awkward phrasing, and 2 indicates the sentence is fluent and almost perfect. 
Provide your rating using this exact format: ``Rating: [[score]]''. \\ \\
{[Sentence Fragment Start]} \\
\textcolor{gray}{[Sentence goes here]} \\ 
{[Sentence Fragment End]} \\
\end{tcolorbox}

\clearpage